\begin{corollary}[Lee--Yang 核的 Stokes 全息：零点由边界决定]\label{cor:xi-hypercube-leyang-stokes-holography}
\leanverified{paper\_xi\_hypercube\_leyang\_stokes\_holography}
在定理 \ref{thm:xi-hypercube-walsh-joukowsky-leyang-explicit} 的设定下，若权函数取为集合 \(B\subseteq\mathcal H_k\) 的指示函数 \(w=\ind_B\)，则对每个 \(x\in\mathcal H_k\)，函数 \(u\mapsto Z_x(u)\) 的 Taylor 系数（等价地 \(t\mapsto P_x(t)\) 的全部系数）可由 \(B\) 的有向边界 \(\partial\overrightarrow{B}\) 的通量数据唯一确定（定理 \ref{thm:xi-hypercube-fourier-walsh-boundary-parseval}）。
因此，对每个固定的 \(x\)，Lee--Yang 零点集合 \(\{u:\ Z_x(u)=0\}\) 亦由 \(\partial\overrightarrow{B}\) 的边界通量唯一确定。
\end{corollary}
\begin{proof}
当 \(w=\ind_B\) 时，
\[
\widehat w(S)=\sum_{a\in\mathcal H_k}\ind_B(a)\chi_S(a)=\sum_{a\in B}\chi_S(a)
=2^k\,\widehat{\ind_B}(S),
\]
其中 \(\widehat{\ind_B}(S)\) 为定理 \ref{thm:xi-hypercube-fourier-walsh-boundary-parseval} 的归一化 Walsh 系数。
取权向量 \(\mathbf 1=(1,\dots,1)\)，则对任意非空 \(S\subseteq[k]\) 有 \(\sum_{i\in S}1\neq 0\)，从而由定理 \ref{thm:xi-hypercube-fourier-walsh-boundary-parseval} 得到
\(
\widehat w(S)=\sum_{(\omega,\nu)\in\partial\overrightarrow{B}}\alpha_{S,\mathbf 1}(\omega,\nu)
\)，即每个非平凡 Walsh 系数为边界通量。
将其代入定理 \ref{thm:xi-hypercube-walsh-joukowsky-leyang-explicit} 的展开式，得到对每个 \(x\) 的 \(Z_x(u)\)（从而 \(P_x(t)\)）完全由边界通量决定。
解析函数由其系数唯一确定，故其零点集合亦由边界通量唯一确定。证毕。
\end{proof}

\begin{theorem}[仿射对合不变权下的 Lee--Yang 多项式二层分解]\label{thm:xi-hypercube-leyang-affine-orbifold-two-layer}
\leanverified{paper\_xi\_hypercube\_leyang\_affine\_orbifold\_two\_layer}
沿用定理 \ref{thm:xi-hypercube-walsh-joukowsky-leyang-explicit} 的设定，取 \(1\le i<j\le k\)，并在 \(\mathcal H_k\) 上定义仿射对合
\[
\sigma_{i,j}(y)_i:=1-y_j,\qquad
\sigma_{i,j}(y)_j:=1-y_i,\qquad
\sigma_{i,j}(y)_r:=y_r\ (r\neq i,j).
\]
若权函数满足 \(\sigma_{i,j}\)-不变性 \(w\circ\sigma_{i,j}=w\)，则对任意 \(x\in\mathcal H_k\) 与 \(t\in\CC\)（其中 \(t=(1-u)/(1+u)\)），多项式
\[
P_x(t)=\sum_{S\subseteq[k]}\widehat w(S)\,t^{|S|}\chi_S(x)
\]
具有严格分解
\[
P_x(t)=P_x^{\mathrm{even}}(t)+\bigl(1-\chi_{\{i,j\}}(x)\bigr)\,Q_x(t),
\]
其中
\[
P_x^{\mathrm{even}}(t):=\sum_{\substack{S\subseteq[k]\\ |S\cap\{i,j\}|\ \mathrm{even}}}\widehat w(S)\,t^{|S|}\chi_S(x),
\]
以及
\[
Q_x(t):=-\sum_{J\subseteq[k]\setminus\{i,j\}}
\widehat w(J\cup\{i\})\,t^{|J|+1}\,\chi_J(x)\,\chi_{\{j\}}(x).
\]
特别地，若 \(x_i=x_j\)（等价于 \(\chi_{\{i,j\}}(x)=1\)），则 \(P_x(t)=P_x^{\mathrm{even}}(t)\)；
若 \(x_i\neq x_j\)（等价于 \(\chi_{\{i,j\}}(x)=-1\)），则 \(P_x(t)=P_x^{\mathrm{even}}(t)+2Q_x(t)\)。
\end{theorem}
\begin{proof}
记 \(R=[k]\setminus\{i,j\}\)。
由 \(w\circ\sigma_{i,j}=w\) 及变元替换 \(a=\sigma_{i,j}(b)\) 得
\[
\widehat w(J\cup\{i\})
=\sum_{a\in\mathcal H_k}w(a)\chi_{J\cup\{i\}}(a)
=\sum_{b\in\mathcal H_k}w(b)\chi_{J\cup\{i\}}(\sigma_{i,j}b)
=-\sum_{b\in\mathcal H_k}w(b)\chi_{J\cup\{j\}}(b)
=-\widehat w(J\cup\{j\}),
\]
其中第三步使用了恒等式 \(\chi_{J\cup\{i\}}\circ\sigma_{i,j}=-\chi_{J\cup\{j\}}\)。
\par
将 \(P_x(t)\) 按 \(|S\cap\{i,j\}|\) 的奇偶分组。
偶层即给出 \(P_x^{\mathrm{even}}(t)\)。
对奇层，配对 \(S=J\cup\{i\}\) 与 \(S'=J\cup\{j\}\)（\(J\subseteq R\)）并用 \(\widehat w(S')=-\widehat w(S)\)，得到
\[
\widehat w(S)t^{|S|}\chi_S(x)+\widehat w(S')t^{|S'|}\chi_{S'}(x)
=\widehat w(J\cup\{i\})\,t^{|J|+1}\,\chi_J(x)\bigl(\chi_{\{i\}}(x)-\chi_{\{j\}}(x)\bigr).
\]
注意
\(
\chi_{\{i\}}(x)-\chi_{\{j\}}(x)=-(1-\chi_{\{i,j\}}(x))\,\chi_{\{j\}}(x)
\)，
将其代回并对 \(J\) 求和即得奇层可因式分解为 \((1-\chi_{\{i,j\}}(x))Q_x(t)\)，从而得到所示分解。
两种取值情形由 \(\chi_{\{i,j\}}(x)\in\{\pm 1\}\) 立得。证毕。
\end{proof}

\endinput
